\lesson{07/10}
\section{Constructing a Lagrangian}
The guiding principle of this process is the Occam's Razor: nature will realize
the simplest Lorentz-invariant terms in the Lagrangian.

The fundamental quantity to build a Lagrangian theory is \textbf{the action}, defined
before as
$$S = \int\dd^4x\mathcal{L}(\phi,\pd_\mu\phi):$$
in classical mechanics the physical trajectory is the one on which the action (as a 
functional) is minimal; in QM, even if no trajectory exists, this process can be still
used in path integral formulation.

As before stated, the Lagrangian density must be \textit{local}, to guarantee
causality, \textit{real}, to preserve probability, \textit{Poincaré 
invariant}, to ensure that the theory is relativistic and must contain \textit{at
most a second derivative}, again to guarantee causality.

Other invariances are admitted (and usually present) and are called \textit{internal
symmetries}: for example the invariance of a spinor under a global phase shift
guarantees \textbf{charge conservation}.
\COMMENT{
\textbf{Example: The Return of Klein-Gordon}\footnote{probably useless,
maybe take it out}\\
The K-G equation can be interpreted as an eigenvalue equation for a field:
$$(\Box+m^2)\phi(x)=0$$
where the d'Alembert operator $\Box = -P_\mu P^\mu$ is a Casimir operator of
eigenvalue $m^2$.
}
\section{Least action}
The action in natural units is dimensionless, so the Lagrangian density must have
units [m$^{-4}$]=[kg$^4$].

Consider an internal symmetry (i.e. a transformation that does not change the
spacetime coordinates)\footnote{this is equivalent to the small variations of the
trajectory in Lagrangian Mechanics.}:
$$\phi(x)\rightarrow \phi'(x')=\phi(x)+\delta\phi(x),$$
the action
$$S=\int_{}^{}\dd^4x\mathcal{L}$$
is minimal if
$$\delta S =0 = \int_{}^{}\dd^4x\delta\mathcal{L} = 
\int_{}^{}\dd^4x \left[\pdv[]{\mathcal{L}}{\phi}\delta\phi + \pdv{\mathcal{L}}
{(\pd_\mu\phi)}\pd_\mu(\delta\phi)\right],
$$
where the relation $\delta(\pd_\mu\phi) = \pd_\mu(\delta\phi)$ was used;
consider
\begin{gather*}
\pd_\mu\left(\pdv[]{\mathcal{L}}{(\pd_\mu\phi)}\delta\phi\right) = 
\pdv[]{\mathcal{L}}{(\pd_\mu\phi)}\pd_\mu(\delta\phi)+\pd_\mu\left(
\pdv[]{\mathcal{L}}{(\pd_\mu\phi)}\right)\delta\phi \rightarrow\\
\rightarrow\pdv[]{\mathcal{L}}{(\pd_\mu\phi)}\pd_\mu(\delta\phi) = \pd_\mu\left(
\pdv[]{\mathcal{L}}{(\pd_\mu\phi)}\delta\phi
\right) - \pd_\mu\left(\pdv[]{\mathcal{L}}{(\pd_\mu\phi)}\right)\delta\phi:
\end{gather*}
substituting it into the $\delta S=0$ condition yields
$$0=\int_{}^{}\dd^4x\left[\pdv[]{\mathcal{L}}{\phi}-\pd_\mu\pdv[]{\mathcal{L}}
{(\pd_\mu\phi)}\right]\delta\phi + \int_{}^{}\dd^4x\pd_\mu\left(\pdv[]{\mathcal{L}}
{(\pd_\mu\phi)}\delta \phi\right).$$
If the variation of the fields $\delta\phi$ goes to zero at the boundary of the
surface enclosing the dominion of integration (i.e. infinity)\footnote{
this is the case for all physical fields.}, the second integral
is zero by the divergence theorem and the action is stationary if and only if
$$\boxed{\pdv[]{\mathcal{L}}{\phi}-\pd_\mu\pdv[]{L}{(\pd_\mu\phi)}=0}$$
which is called the \textbf{Euler-Lagrange equation} for fields.

\textbf{Note}: the E-L equation is invariant under transformations in the form
$$\mathcal{L}'=\mathcal{L}+\pd_\mu\Lambda^\mu$$
where $\Lambda^\mu$ is any 4-vector function: these are called \textit{canonical
transformations}. A trivial example of these transformations are simple translations
of $\mathcal{L}\rightarrow\mathcal{L}+C$.

\section{Noether's Theorem}
This is one of the pivotal theorems of QFT and all of theoretical physics and it
states

\noindent\fbox{%
    \parbox{\textwidth}{%
    \textit{To each symmetry of a system corresponds a 4-current satisfying a
    continuity equation and a conserved quantity, called Noether charge.}
    }%
}

Consider a transformation that acts both on the coordinates and a field
$$\begin{dcases*}
  x'^\mu = x^\mu + \delta x^\mu\\
  \phi'(x')=\phi(x)+\Delta\phi(x)
  \end{dcases*}$$
which does not change the action $S$\footnote{This is the definition of symmetry}.
Staring with $\phi$ as a solution to the E-L equation, the action is invariant under
the transformation if and only if
$$\Delta S =0=\int_{}^{}\dd^4x\Delta\mathcal{L}+\int_{}^{}\delta(\dd^4x)\mathcal{L}$$
where
$$\Delta\mathcal{L} = \underbrace{\pdv[]{\mathcal{L}}{\phi}\delta\phi+
\pdv[]{\mathcal{L}}{\pd_\mu\phi}\pd_\mu\delta\phi}_{
= \delta\mathcal{L}} + \pd_\mu\mathcal{L}\delta x^\mu= \underbrace{\left(\pdv[]{
\mathcal{L}}{\phi} - \pd_\mu\pdv{\mathcal{L}}{\pd_\mu\phi}\right)}_{=0\text{ E-L}}
\delta\phi+\pd_\mu\left(\pdv[]{\mathcal{L}}{\pd_\mu\phi}\delta\phi\right)+
(\pd_\mu\mathcal{L})\delta x^\mu.$$
In this case, since the transformation acts on the coordinates, the variation of
the integration measure must be considered:
$$\delta(\dd^4x)= \dd^4x'-\dd^4x$$
using the Jacobian of the coordinate transformation
$$\dd^4x' = |\det(J(x))|\dd^4x = \left|\det\pdv{x'^\mu}{x^\nu}\right|\dd^4x$$
$$
\delta(\dd^4x)=\dd^4x'-\dd^4x\approx \pd_\mu\delta x^\mu\dd^4x,
$$
since
$$\pdv{x'^\mu}{x^\mu} = \pdv{(x^\mu+\delta x^\mu)}{x^\nu} = g\indices{^\mu_\nu}+
\pd_\nu\delta x^\mu,$$
the property $\det A=e^{\Tr(\ln A)}$ yields
$$\det J(x) \underset{\delta x^\mu\to0}\approx\exp\left[\Tr\pd_\nu\delta x^\mu\right]
\sim 1+\pd_\mu\delta x^\mu.$$
From this follows that the variation of the measure is
$$\delta(\dd^4x) \approx \pd_\mu\delta x^\mu\dd^4x.$$

So the action is stationary if and only if
\begin{gather*}
0 = \Delta S = \int_{}^{}\dd^4x \left[\pd_\mu\left(\pdv[]{\mathcal{L}}{\pd_\mu\phi}
\delta\phi\right)+\underbrace{\pd_\mu\mathcal{L}\delta x^\mu + \mathcal{L}\pd_\mu
\delta x^\mu}_{=\pd_\mu(\mathcal{L}\delta x^\mu)}\right] = \\
= \int_{}^{}\dd^4x \pd_\mu\left(\pdv[]{\mathcal{L}}{\pd_\mu\phi}\delta\phi+
\mathcal{L}\delta x^\mu\right) = 0
\end{gather*}
this happens if and only if
$$\pd_\mu\left(\pdv[]{\mathcal{L}}{\pd_\mu\phi}\delta\phi+
\mathcal{L}\delta x^\mu\right)= 0. $$
The function in parenthesis is \textit{the Noether current}
$$J^\mu =\pdv[]{\mathcal{L}}{\pd_\mu\phi}\delta\phi+\mathcal{L}\delta x^\mu $$
and the found relation is its continuity equation, as
$$\pd_\mu J^\mu = \pd_0 J^0 + \pd_i J^i = \pdv[]{\rho}{t}+\divergence\mathbf{j}=0,$$
where $J^0:=\rho$, the density associated with the 3-current
$\mathbf{j}$.

Integrating the continuity equation yields
$$0=\int_{}^{}\dd^3x \pd_\mu J^\mu = \int_{}^{}\dd^3x \pd_0J^0+
\underbrace{\int_{}^{}\dd^3x\pd_i J^i}_{=0\text{ by div thm}}=
\int_{}^{}d^3x\pdv[]{J^0}{t} = \dv[]{}{t}\int_{}^{}\dd^3xJ^0:=\dv[]{Q}{t}=0$$
so the Noether charge which is conserved is
$$Q = \int_{}^{}d^3xJ^0(x^\mu).$$
It is useful to rewrite the Noether current in terms of $\Delta\phi$:
recalling that
$$\delta\phi = \Delta\phi - \pd_\nu\phi\delta x^\nu$$
the current becomes
$$J^\mu = \underbrace{\left(\mathcal{L}g\indices{^\mu_\nu}-\pdv[]{\mathcal{L}}{
\pd_\mu\phi}\pd_\nu\phi\right)}_{:=-T\indices{^\mu_\nu}}\delta x^\nu +
\pdv[]{\mathcal{L}}{\pd_\mu\phi}\Delta\phi$$
where $T\indices{^\mu_\nu}$ is called the \textit{stress energy tensor}; the sign
is a convention that ensures $T\indices{^0_0}\geq0$.
